% संपूर्ण मराठी ग्रंथातील प्रकरण 8: सिद्धता-पद्धती
% हा संपादनयोग्य स्रोतखंड आहे. संकलनासाठी संपूर्ण स्रोतसंग्रहातील
% अक्षररूपे, पूर्वभाग, चिन्हव्याख्या आणि आकृती-स्रोत आवश्यक आहेत.
% मूळ: Open Logic Project; मजकूर आणि रूपांतर CC BY 4.0.
% यंत्रानुवाद, दुरुस्ती आणि तपासणी: OpenAI Codex —
% GPT-5.6 Sol आणि GPT-6 Sol, दोन्ही Ultra effort.
% दुरुस्ती-पुनरावलोकन: GPT-6.1 Sol, Ultra effort.
\chapter{सिद्धता-पद्धती}\label{fol:prf::chap}
\begin{editorial}
या प्रकरणात निष्पत्ती-पद्धतींविषयीची सर्वसाधारण सामग्री
एकत्रित केली आहे. एखादी विशिष्ट पद्धत वापरणारे पाठ्यपुस्तक,
ज्या प्रकरणात त्या पद्धतीचा परिचय करून दिला आहे त्या प्रकरणाच्या
सुरुवातीस प्रस्तावना-विभाग आणि संबंधित आढावा-विभाग समाविष्ट करू शकते.
\end{editorial}
\section{परिचय}\label{fol:prf:int:sec}
तर्कप्रणालींमध्ये सामान्यतः चिन्हार्थमीमांसा आणि निष्पत्ती-पद्धती हे दोन्ही घटक असतात. चिन्हार्थमीमांसेत सत्यता, पूर्ततायोग्यता,
वैधता आणि तार्किक निष्पन्नता अशा संकल्पनांचा विचार केला जातो.
तार्किक निष्पन्नता आणि वैधता प्रस्थापित करण्यासाठी निव्वळ विन्यासात्मक
रीत उपलब्ध करून देणे हा निष्पत्ती-पद्धतींचा उद्देश असतो. त्या
निव्वळ विन्यासात्मक असतात याचा अर्थ असा की अशा पद्धतीतील
निष्पत्ती ही एक सांत विन्यासात्मक वस्तू असते; सहसा ती
वाक्ये किंवा सूत्रे यांचा अनुक्रम (किंवा इतर एखादी सांत
मांडणी) असते. वाक्ये किंवा सूत्रे यांचा कोणताही दिलेला
अनुक्रम किंवा मांडणी ``योग्य'' आहे, हे यांत्रिक रीतीने पडताळता येते,
हा चांगल्या निष्पत्ती-पद्धतींचा गुणधर्म असतो.
प्रथम-क्रम तर्कशास्त्रासाठीच्या सर्वांत सोप्या (आणि इतिहासात प्रथम
आलेल्या) निष्पत्ती-पद्धती \emph{स्वयंसिद्धकीय} होत्या. अशा
पद्धतीत सूत्रांचा एखादा अनुक्रम तेव्हाच निष्पत्ती मानला
जातो, जेव्हा त्यातील प्रत्येक स्वतंत्र सूत्र एकतर
``स्वयंसिद्धकांच्या'' एका निश्चित संचात असते किंवा त्या अनुक्रमात
त्याच्या आधी आलेल्या सूत्रांपासून ठरावीक संख्येतील ``निगमन
नियमां''पैकी एखाद्या नियमाने निष्पन्न होते. शिवाय, सूत्र
स्वयंसिद्धक आहे का आणि ते इतर सूत्रांपासून एखाद्या निगमन
नियमानुसार योग्य रीतीने निष्पन्न होते का, हे यांत्रिक रीतीने पडताळता
येते. स्वयंसिद्धकीय निष्पत्ती-पद्धतींचे वर्णन करणे सोपे असते आणि
अधिउपपत्तीच्या पातळीवरही त्या हाताळणे सोपे असते; पण त्यांतील
निष्पत्त्या वाचणे आणि समजून घेणे कठीण असते, तसेच त्या तयार करणेही
कठीण असते.
निष्पत्त्या तयार करणे सोपे व्हावे किंवा पूर्ण झालेल्या
निष्पत्त्या समजून घेणे सोपे व्हावे, या उद्देशाने इतर
निष्पत्ती-पद्धती विकसित करण्यात आल्या आहेत. नैसर्गिक निगमन,
सत्यता-वृक्ष---ज्यांना टॅब्लो सिद्धता असेही म्हणतात---आणि क्रमवर्ती कलन
ही त्यांची उदाहरणे आहेत. काही निष्पत्ती-पद्धती विशेषतः
यांत्रिकीकरण डोळ्यांसमोर ठेवून रचल्या जातात. उदाहरणार्थ, निराकरण पद्धत
संगणकीय प्रणालीत कार्यान्वित करणे सोपे असते (पण तिच्यातील
निष्पत्त्या समजून घेणे जवळजवळ अशक्य असते). यांपैकी बहुतेक इतर
निष्पत्ती-पद्धती निष्पत्त्यांना अनुक्रमांऐवजी सूत्रांचे
वृक्ष म्हणून दर्शवतात. त्यामुळे निष्पत्तीचा कोणता भाग इतर
कोणत्या भागांवर अवलंबून आहे, हे पाहणे सोपे होते.
म्हणून प्रथम-क्रम तर्कशास्त्रासारखी एखादी तर्कप्रणाली घेतल्यास,
वाक्य हे \emph{प्रमेय} असणे म्हणजे काय आणि वाक्य हे
इतर काही वाक्यांपासून निष्पन्न करता येण्याजोगे असणे म्हणजे काय, यांची
वेगवेगळ्या निष्पत्ती-पद्धती वेगवेगळी नेमकी स्पष्टीकरणे देतील. हे
कसेही केले (स्वयंसिद्धकीय निष्पत्त्या, नैसर्गिक निगमने, क्रमवर्ती
निष्पत्त्या, सत्यता-वृक्ष किंवा निराकरण-खंडने यांच्या साहाय्याने),
तरी हे संबंध वैधता आणि तार्किक निष्पन्नता या चिन्हार्थविषयक संकल्पनांशी
जुळावेत, अशी आपली अपेक्षा असते. ``$\varphi$~हे प्रमेय आहे'' यासाठी
$\Proves \varphi$ आणि ``$\varphi$~हे~$\Gamma$ पासून निष्पन्न करता येण्याजोगे आहे'' यासाठी
``$\Gamma \Proves \varphi$'' लिहू या. $\Proves$~ची व्याख्या कोणतीही असो,
ती $\Entails$~शी जुळावी, म्हणजे:
\begin{enumerate}
\item $\Proves \varphi$ जर आणि फक्त जर $\Entails \varphi$
\item $\Gamma \Proves \varphi$ जर आणि फक्त जर $\Gamma \Entails \varphi$
\end{enumerate}
वरील विधानांतील ``फक्त जर'' दिशेला \emph{निर्दोषता} म्हणतात.
निष्पत्ती-पद्धत तेव्हा निर्दोष असते, जेव्हा निष्पन्नतेमुळे
तार्किक निष्पन्नतेची (किंवा वैधतेची) हमी मिळते. प्रत्येक समाधानकारक
निष्पत्ती-पद्धत निर्दोष असलीच पाहिजे; सदोष निष्पत्ती-पद्धती
मुळीच उपयोगी नसतात. कारण निष्पत्तीचा संपूर्ण उद्देशच वैधतेची
किंवा तार्किक निष्पन्नतेची विन्यासात्मक हमी देणे हा असतो. आपण मांडत
असलेल्या निष्पत्ती-पद्धतींची निर्दोषता पुढे सिद्ध करू.
उलट ``जर'' दिशेलाही महत्त्व आहे: तिला \emph{संपूर्णता} म्हणतात.
एखादी संपूर्ण निष्पत्ती-पद्धत इतकी समर्थ असते की $\varphi$~वैध असेल
तेव्हा $\varphi$~हे प्रमेय आहे, तसेच $\Gamma \Entails \varphi$ असेल तेव्हा
$\Gamma \Proves \varphi$, हे ती दाखवू शकते. संपूर्णता प्रस्थापित करणे अधिक
कठीण असते आणि काही तर्कप्रणालींसाठी संपूर्ण निष्पत्ती-पद्धतीच
नसतात. प्रथम-क्रम तर्कशास्त्रासाठी मात्र त्या आहेत. कुर्ट गोडेल यांनी
आपल्या १९२९ सालच्या प्रबंधात प्रथम-क्रम तर्कशास्त्राच्या
निष्पत्ती-पद्धतीची संपूर्णता सर्वप्रथम सिद्ध केली.
निष्पत्ती-पद्धतींशी जोडलेली आणखी एक संकल्पना म्हणजे
\emph{सुसंगतता}. वाक्यांच्या एखाद्या संचापासून कोणतीही गोष्ट
निष्पन्न करता येत असेल तर तो संच विसंगत म्हणतात; अन्यथा तो सुसंगत
असतो. विसंगतता ही अपूर्ततायोग्यतेची विन्यासात्मक समकक्ष संकल्पना आहे:
जसे अपूर्ततायोग्य संच चांगल्या उपपत्ती ठरत नाहीत, तसेच वाक्यांचे
विसंगत संचही चांगल्या उपपत्ती ठरत नाहीत; त्यांच्यात मूलभूत स्वरूपाचा
दोष असतो. वाक्यांचे सुसंगत संच सत्य किंवा उपयोगी असतीलच असे
नाही, पण ते निदान तार्किक उपयोगितेचा किमान उंबरठा ओलांडतात. वेगवेगळ्या
निष्पत्ती-पद्धतींसाठी वाक्यांच्या संचांच्या सुसंगततेची नेमकी
व्याख्या वेगळी असू शकते; पण~$\Proves$ प्रमाणेच सुसंगतताही तिच्या
चिन्हार्थविषयक समकक्ष संकल्पनेशी, म्हणजे पूर्ततायोग्यतेशी, जुळली पाहिजे.
$\Gamma$ सुसंगत आहे जर आणि फक्त जर तो पूर्ततायोग्य आहे, हे नेहमी खरे
असावे अशी आपली अपेक्षा आहे. येथे ``फक्त जर'' दिशा संपूर्णतेइतकीच ठरते
(सुसंगततेमुळे पूर्ततायोग्यतेची हमी मिळते), आणि ``जर'' दिशा
निर्दोषतेइतकीच ठरते (पूर्ततायोग्यतेमुळे सुसंगततेची हमी मिळते).
प्रत्यक्षात, अभिजात प्रथम-क्रम तर्कशास्त्रासाठी निर्दोषता आणि संपूर्णता
यांची ही दोन रूपे सममूल्य आहेत.
\section{क्रमवर्ती कलन}\label{fol:prf:seq:sec}
अनेक निष्पत्ती-पद्धती वाक्यांच्या मांडण्यांवर कार्य करतात,
तर क्रमवर्ती कलन \emph{क्रमवर्तींवर} कार्य करते. क्रमवर्ती ही पुढील रूपाची
अभिव्यक्ती असते:
\[
\varphi_1, \dots, \varphi_m \Sequent \psi_1, \dots, \psi_m,
\]
म्हणजे वाक्यांच्या अशा दोन अनुक्रमांची जोडी, ज्यांना क्रमवर्ती
चिन्ह~$\Sequent$ वेगळे करते. यांपैकी कोणताही अनुक्रम रिकामा असू शकतो.
क्रमवर्ती कलनातील निष्पत्ती म्हणजे क्रमवर्तींचा वृक्ष होय. त्यातील
सर्वांत वरच्या क्रमवर्ती विशेष रूपाच्या असतात (त्यांना ``आरंभीच्या
क्रमवर्ती'' किंवा ``स्वयंसिद्धके'' म्हणतात), आणि इतर प्रत्येक क्रमवर्ती
तिच्या लगत वर असलेल्या क्रमवर्तींपासून निगमन नियमांपैकी एखाद्या नियमाने
निष्पन्न होते. निगमन नियम एकतर क्रमवर्तींतील वाक्ये हाताळतात
(डाव्या किंवा उजव्या बाजूला त्यांची भर घालतात, त्यांना काढतात किंवा त्यांचा क्रम
बदलतात), किंवा नियमाच्या निष्कर्षात एखादे जटिल सूत्र आणतात.
उदाहरणार्थ, $\LeftR{\land}$ नियमामुळे $\varphi, \Gamma \Sequent \Delta$ पासून
$\varphi \land \psi, \Gamma \Sequent \Delta$ हा निष्कर्ष काढता येतो; आणि
$\RightR{\lif}$ नियमामुळे $\varphi, \Gamma \Sequent \Delta, \psi$ पासून
$\Gamma \Sequent \Delta, \varphi \lif \psi$ हा निष्कर्ष काढता येतो. हे कोणत्याही
$\Gamma$, $\Delta$, $\varphi$ आणि~$\psi$ साठी लागू आहे. (विशेषतः, $\Gamma$
आणि~$\Delta$ रिकामे असू शकतात.)

क्रमवर्ती कलनावर आधारित $\Proves$ संबंधाची व्याख्या पुढीलप्रमाणे केली
जाते: $\Gamma \Proves \varphi$ तेव्हा आणि केवळ तेव्हाच, जेव्हा असा एखादा
अनुक्रम~$\Gamma_0$ अस्तित्वात असतो की $\Gamma_0$ मधील प्रत्येक $\varphi$
हा~$\Gamma$ मध्ये असतो आणि ज्याच्या मुळाशी $\Gamma_0 \Sequent \varphi$ ही
क्रमवर्ती आहे अशी निष्पत्ती अस्तित्वात असते. $\Sequent \varphi$ या क्रमवर्तीची
निष्पत्ती असेल तर क्रमवर्ती कलनात $\varphi$ हे प्रमेय असते. उदाहरणार्थ,
$\Proves (\varphi \land \psi) \lif \varphi$ हे दाखवणारी निष्पत्ती पुढे दिली आहे:
\begin{prooftree}
  \Axiom$\varphi \fCenter \varphi$
  \RightLabel{\LeftR{\land}}
  \UnaryInf$\varphi \land \psi \fCenter \varphi$
  \RightLabel{\RightR{\lif}}
  \UnaryInf$\fCenter (\varphi \land \psi) \lif \varphi$
\end{prooftree}

प्रत्येक $\varphi \in \Gamma_0$ हे~$\Gamma$ मध्ये असेल अशा एखाद्या
$\Gamma_0 \Sequent$ ची निष्पत्ती अस्तित्वात असेल (म्हणजे
क्रमवर्तीची उजवी बाजू रिकामी असेल), तर क्रमवर्ती कलनात संच~$\Gamma$
विसंगत असतो. \RightR{\Weakening} नियम वापरून विसंगत संचापासून कोणतेही
वाक्य निष्पन्न करता येते.

क्रमवर्ती कलनाची निर्मिती १९३० च्या दशकात गरहार्ड गेंटझेन यांनी केली.
त्याची रचना पद्धतशीर आणि सममित असल्यामुळे निष्पत्त्यांची उपपत्ती
विकसित करण्यासाठी ते अतिशय उपयुक्त आकारिक साधन आहे. क्रमवर्ती कलनात
निष्पत्त्या शोधणे तुलनेने सोपे असते; पण या निष्पत्त्या अनेकदा
वाचायला कठीण असतात आणि त्यांचा सिद्धतांशी असलेला संबंधही काही वेळा
सहज दिसत नाही. तरी निष्पत्ती-पद्धतींकडे पाहण्याची ही अतिशय सुबक
रीत ठरली आहे, आणि अनेक तर्कप्रणालींसाठी क्रमवर्ती-कलन पद्धती उपलब्ध आहेत.









\section{नैसर्गिक निगमन}\label{fol:prf:ntd:sec}

नैसर्गिक निगमन ही प्रत्यक्ष तर्कविचाराचे अनुकरण करण्यासाठी घडवलेली
निष्पत्ती-पद्धत आहे (विशेषतः गणितज्ञ वापरतात त्या नियमबद्ध
तर्कविचाराचे). प्रत्यक्ष तर्कविचार अनेक ``नैसर्गिक'' पद्धतींनी पुढे जातो.
उदाहरणार्थ, प्रकरणांनुसार सिद्धतेमुळे विकल्पयोगी आधारविधानाच्या आधारे
निष्कर्ष प्रस्थापित करता येतो: विकल्पयोगाच्या प्रत्येक घटकापासून
तो निष्कर्ष निष्पन्न होतो, हे स्वतंत्रपणे दाखवले जाते. अप्रत्यक्ष सिद्धतेत
निष्कर्षाचे नकरण व्याघाताकडे नेते हे दाखवून निष्कर्ष प्रस्थापित करता येतो.
सोपाधिक सिद्धतेत पूर्वांगापासून उत्तरांग निष्पन्न होते हे दाखवून
``जर \dots तर \dots'' असा सोपाधिक दावा प्रस्थापित केला जातो. नैसर्गिक
निगमन म्हणजे अशा नैसर्गिक अनुमानांपैकी काहींचे आकारिकीकरण होय. प्रत्येक
तार्किक संयोजक आणि संख्यापक यांना दोन नियम असतात—एक प्रवेशन नियम आणि एक
विलोपन नियम—आणि हे नियम प्रत्येकी अशाच एका नैसर्गिक अनुमान-पद्धतीशी
जुळतात. उदाहरणार्थ, $\Intro{\lif}$ सोपाधिक सिद्धतेशी, तर $\Elim{\lor}$
प्रकरणांनुसार सिद्धतेशी जुळतो. $\Elim{\land}$ हा विशेषतः सोपा नियम आहे;
त्यामुळे $\varphi \land \psi$ पासून~$\varphi$ (किंवा $\psi$) असा निष्कर्ष काढता येतो.

नैसर्गिक निगमनाला इतर निष्पत्ती-पद्धतींपासून वेगळे करणारे एक
वैशिष्ट्य म्हणजे त्यात गृहीतकांचा होणारा वापर. नैसर्गिक निगमनातील निष्पत्ती म्हणजे सूत्रांचा वृक्ष होय. या सूत्रांच्या
वृक्षाच्या मुळाशी एकच सूत्र असते, आणि वृक्षाची ``पाने'' ही अशी
सूत्रे असतात की ज्यांच्यापासून निष्कर्ष काढला जातो. नैसर्गिक
निगमनात पानांवरील काही सूत्रे ही निष्पत्तीमध्ये भूमिका बजावतात;
पण निष्पत्ती निष्कर्षापर्यंत पोहोचेपर्यंत ती ``वापरून संपवलेली'' असतात.
प्रत्यक्ष तर्कविचारात फक्त थोड्या काळापुरती लागू असणारी गृहीतके मांडण्याच्या
प्रथेशी हे जुळते. उदाहरणार्थ, प्रकरणांनुसार सिद्धतेत आपण विकल्पयोगाच्या
प्रत्येक घटकाची सत्यता गृहीत धरतो; सोपाधिक सिद्धतेत पूर्वांगाची सत्यता
गृहीत धरतो; आणि अप्रत्यक्ष सिद्धतेत निष्कर्षाच्या नकरणाची सत्यता गृहीत
धरतो. अशी तात्पुरती गृहीतके मांडणे आणि एखादी मधली पायरी प्रस्थापित करण्यासाठी
ती नंतर बाजूला काढणे, हे नैसर्गिक निगमनाचे खास वैशिष्ट्य आहे. नैसर्गिक
निगमनातील निष्पत्तीच्या पानांवरील सूत्रांना गृहीतके म्हणतात,
आणि काही निगमन नियम त्यांना ``मुक्त'' करू शकतात. उदाहरणार्थ, काही
गृहीतकांपासून~$\psi$ ची अशी निष्पत्ती असेल आणि त्या गृहीतकांत~$\varphi$
असेल, तर $\Intro{\lif}$ नियमामुळे~$\varphi \lif \psi$ असा निष्कर्ष काढता येतो
आणि~$\varphi$ या रूपातील कोणतेही गृहीतक मुक्त करता येते. (कोणती गृहीतके
कोणत्या अनुमानाच्या वेळी मुक्त केली जातात याची नोंद ठेवण्यासाठी, त्या
अनुमानाला आणि त्याने मुक्त केलेल्या गृहीतकांना
एक क्रमांक देतो.) निष्पत्तीच्या शेवटी जी गृहीतके मुक्त न केलेली
राहतात, ती निष्कर्षाच्या सत्यतेसाठी एकत्रितपणे पुरेशी असतात; म्हणून
निष्पत्ती ही प्रस्थापित करते की तिची मुक्त न केलेली गृहीतके तिचा निष्कर्ष तार्किकरीत्या निष्पन्न करतात.

नैसर्गिक निगमनावर आधारित $\Gamma \Proves \varphi$ हा संबंध तेव्हा आणि केवळ
तेव्हाच लागू होतो, जेव्हा अशी निष्पत्ती असते की त्यात वृक्षातील
शेवटचे वाक्य हे~$\varphi$ असते आणि मुक्त न केलेले प्रत्येक पान
हे~$\Gamma$ मध्ये असते. नैसर्गिक निगमनात $\varphi$ हे तेव्हा आणि केवळ तेव्हाच
प्रमेय असते, जेव्हा अशी निष्पत्ती असते की त्यातील शेवटचे
वाक्य हे~$\varphi$ असते आणि सर्व गृहीतके मुक्त झालेली असतात.
उदाहरणार्थ, $\Proves (\varphi \land \psi) \lif \varphi$ हे दाखवणारी निष्पत्ती
पुढे दिली आहे:
\begin{prooftree}
  \AxiomC{$\Discharge{\varphi \land \psi}{1}$}
  \RightLabel{\Elim{\land}}
  \UnaryInfC{$\varphi$}
  \DischargeRule{\Intro{\lif}}{1}
  \UnaryInfC{$(\varphi \land \psi) \lif \varphi$}
\end{prooftree}
$1$ ही खूण दाखवते की $\varphi \land \psi$ हे गृहीतक \Intro{\lif} अनुमानाच्या
वेळी मुक्त केले जाते.

नैसर्गिक निगमनात संच~$\Gamma$ तेव्हा आणि केवळ तेव्हाच विसंगत असतो,
जेव्हा $\Gamma \Proves \lfalse$. \FalseInt{} नियमामुळे विसंगत संचापासून
कोणतेही वाक्य निष्पन्न करता येते.

नैसर्गिक निगमन पद्धती १९३० च्या दशकात गरहार्ड गेंटझेन आणि स्टॅनिस्लॉ
यास्कोव्हस्की (Stanis\l{}aw Ja\'skowski) यांनी विकसित केल्या; पुढे डॅग
प्रावित्झ आणि फ्रेडरिक फिच यांनी त्यांचा आणखी विकास केला. त्यांतील अनुमाने
सिद्धतेच्या नैसर्गिक पद्धतींचे अनुकरण करत असल्यामुळे तत्त्वज्ञ नैसर्गिक
निगमनाला पसंती देतात.
फिच यांनी विकसित केलेल्या आवृत्त्या तर्कशास्त्राच्या प्रास्ताविक
पाठ्यपुस्तकांमध्ये अनेकदा वापरल्या जातात. तर्कशास्त्राच्या तत्त्वज्ञानात
नैसर्गिक निगमनाचे नियम तार्किक कारकांचे अर्थ देतात, असे काही वेळा मानले
गेले आहे (``सिद्धता-उपपत्तीय चिन्हार्थमीमांसा'').









\section{टॅब्लो}\label{fol:prf:tab:sec}

अनेक निष्पत्ती-पद्धती वाक्यांच्या मांडणीवर काम करतात, तर
टॅब्लो चिन्हांकित सूत्रांवर काम करतात. चिन्हांकित सूत्र म्हणजे
सत्यतामूल्याचे चिन्ह ($\True$ किंवा $\False$) आणि वाक्य यांची जोडी:
\[\sFmla{\True}{\varphi} \text{ किंवा } \sFmla{\False}{\varphi}.
\]
टॅब्लोमध्ये चिन्हांकित सूत्रे खालीच्या दिशेने शाखा फुटणाऱ्या
वृक्षात मांडलेली असतात. त्याची सुरुवात काही \emph{गृहीतकांपासून} होते आणि
त्यांच्या वर असलेल्या चिन्हांकित सूत्रांपैकी एखाद्याला निगमन नियमांपैकी एक
नियम लावल्यामुळे मिळणाऱ्या चिन्हांकित सूत्रांनी तो पुढे चालू राहतो.
प्रत्येक नियमामुळे एखाद्या शाखेच्या शेवटी एक किंवा अधिक चिन्हांकित सूत्रे
जोडता येतात, किंवा दोन चिन्हांकित सूत्रे शेजारी शेजारी जोडता येतात---दुसऱ्या
स्थितीत शाखा दोन शाखांत फुटते आणि नव्याने जोडलेली दोन चिन्हांकित सूत्रे त्या दोन शाखांची टोके बनतात.

संमिश्र चिन्हांकित सूत्राला नियम लावल्यावर तात्काळ उप-सूत्रे असलेली
चिन्हांकित सूत्रे जोडली जातात. हे नियम जोड्यांनी येतात—दोन्ही चिन्हांसाठी
प्रत्येकी एक नियम. उदाहरणार्थ, $\TRule{\True}{\land}$ हा नियम
$\sFmla{\True}{\varphi \land \psi}$ ला लागू होतो आणि
$\sFmla{\True}{\varphi \land \psi}$ असलेल्या कोणत्याही शाखेच्या शेवटी
$\sFmla{\True}{\varphi}$ आणि~$\sFmla{\True}{\psi}$ ही दोन्ही चिन्हांकित सूत्रे
जोडू देतो; तसेच $\TRule{\False}{\varphi \land \psi}$ हा नियम
$\sFmla{\False}{\varphi}$ आणि $\sFmla{\False}{\psi}$ शेजारी शेजारी जोडून शाखा
फोडू देतो. टॅब्लोच्या प्रत्येक शाखेत $\sFmla{\True}{\varphi}$ आणि
$\sFmla{\False}{\varphi}$ ही चिन्हांकित सूत्रांची जुळणारी जोडी असेल, तर तो
बंद असतो.

टॅब्लोवर आधारित $\Proves$ संबंधाची व्याख्या पुढीलप्रमाणे केली जाते:
$\Gamma \Proves \varphi$ तेव्हा आणि केवळ तेव्हाच, जेव्हा असा काही सांत संच
~$\Gamma_0 = \{\psi_1, \dots, \psi_n\} \subseteq \Gamma$ असतो की पुढील
गृहीतकांसाठी बंद टॅब्लो असतो:
\[\{\sFmla{\False}{\varphi}, \sFmla{\True}{\psi_1}, \dots, \sFmla{\True}{\psi_n}\}
\]
उदाहरणार्थ, $\Proves (\varphi \land \psi) \lif \varphi$ हे दाखवणारा बंद टॅब्लो
पुढे दिला आहे:
\begin{oltableau}
  [\sFmla{\False}{(\varphi \land \psi) \lif \varphi}, just = \TAss
    [\sFmla{\True}{\varphi \land \psi}, just = {\TRule{\False}{\lif}[1]}
      [\sFmla{\False}{\varphi}, just = {\TRule{\False}{\lif}[1]}
        [\sFmla{\True}{\varphi}, just = {\TRule{\True}{\land}[2]}
          [\sFmla{\True}{\psi}, just = {\TRule{\True}{\land}[2]}, close]
        ]
      ]
    ]
  ]
\end{oltableau}

काही $\psi_i \in \Gamma$ साठी पुढील गृहीतकांचा बंद टॅब्लो असेल, तर
संच $\Gamma$ हा टॅब्लो कलनात विसंगत असतो:
\[\{\sFmla{\True}{\psi_1}, \dots, \sFmla{\True}{\psi_n}\}
\]

टॅब्लोचा शोध १९५० च्या दशकात एव्हर्ट बेथ आणि याक्को हिंटिक्का यांनी
स्वतंत्रपणे लावला; रेमंड स्मलियन यांनी त्यांचे सुलभीकरण करून ते लोकप्रिय
केले. टॅब्लो तयार करणे ही अत्यंत पद्धतशीर प्रक्रिया असल्यामुळे त्यांचा
वापर अगदी सोपा आहे. टॅब्लोच्या पद्धतशीर स्वरूपामुळे त्यांची
संगणकीय अंमलबजावणीही सुलभ होते. तथापि, टॅब्लो वाचणे अनेकदा कठीण
असते आणि त्यांचा सिद्धतांशी असलेला संबंध काही वेळा सहज दिसत नाही. ही
दृष्टीही बरीच व्यापक आहे आणि अनेक भिन्न तर्कशास्त्रांसाठी टॅब्लो
पद्धती उपलब्ध आहेत. दिलेली वाक्ये (किंवा त्यांचे संच) पूर्ण करणाऱ्या
संरचना शोधण्यासही टॅब्लो मदत करतात: संच पूर्ततायोग्य असेल,
तर त्याचा बंद टॅब्लो असणार नाही; म्हणजेच कोणत्याही टॅब्लोमध्ये
एक खुली शाखा असेल. त्या शाखेवर लागू करता येणारा प्रत्येक नियम लागू केलेला
असेल, तर खुल्या शाखेवरून पूर्ति करणारी संरचना ``वाचून काढता'' येते.
क्रमवर्ती कलनाशीही याचा अत्यंत निकटचा संबंध आहे: मूलतः, बंद टॅब्लो
म्हणजे उलटी लिहिलेली क्रमवर्ती कलनातील संक्षेपित निष्पत्ती होय.









\section{स्वयंसिद्धकीय निष्पत्ती}\label{fol:prf:axd:sec}

स्वयंसिद्धकीय निष्पत्त्या या सर्वांत जुन्या आणि सोप्या तार्किक
निष्पत्ती-पद्धती आहेत. त्यांतील निष्पत्त्या म्हणजे केवळ
वाक्यांचे अनुक्रम असतात. वाक्यांचा अनुक्रम तेव्हा योग्य
निष्पत्ती मानला जातो, जेव्हा त्यातील प्रत्येक वाक्य~$\varphi$ पुढील
अटींपैकी एक अट पूर्ण करते:
\begin{enumerate}
\item $\varphi$ हे स्वयंसिद्धक असते, किंवा
\item $\varphi$ हे दिलेल्या वाक्यांच्या संच~$\Gamma$ चा घटक असते, किंवा
\item $\varphi$ ला निगमन नियमामुळे समर्थन मिळते.
\end{enumerate}
$\varphi$ हे स्वयंसिद्धक असण्यासाठी त्याचे रूप काही ठरावीक वाक्य
रूपबंधांपैकी एक असावे लागते. प्रथम-क्रम तर्कशास्त्रासाठी समाधानकारक
(निर्दोष आणि संपूर्ण) निष्पत्ती-पद्धत देणारे स्वयंसिद्धक-रूपबंधांचे
अनेक संच आहेत. त्यांपैकी काही संच, ते ज्या संयोजकांना नियंत्रित करतात
त्यांनुसार मांडलेले असतात; उदाहरणार्थ, पुढील रूपबंध
\[\varphi \lif (\psi \lif \varphi) \qquad \psi \lif (\psi \lor \chi) \qquad (\psi \land \chi) \lif \psi
\]
हे $\lif$, $\lor$ आणि~$\land$ यांना नियंत्रित करणारे सर्वसाधारण स्वयंसिद्धक
आहेत. काही स्वयंसिद्धकीय पद्धतींमध्ये स्वयंसिद्धकांची संख्या कमीत कमी
ठेवण्याचे उद्दिष्ट असते. कोणते संयोजक आदिम म्हणून स्वीकारले आहेत यानुसार,
एकाच स्वयंसिद्धकाची पद्धत मिळणेही शक्य असते.

निगमन नियम म्हणजे असे सोपाधिक विधान की जे निष्पत्तीमधील
वाक्य समर्थित असण्यासाठी पुरेशी अट देते. विध्यात्मक अनुमान हा असा
एक अतिशय प्रचलित नियम आहे: $\varphi$ आणि $\varphi \lif \psi$ यांना आधीच समर्थन मिळाले
असेल, तर $\psi$ लाही समर्थन मिळते, असे तो सांगतो. म्हणजे, निष्पत्तीमध्ये वाक्य~$\psi$ असलेल्या ओळीला समर्थन मिळते, मात्र $\varphi$ आणि
$\varphi \lif \psi$ ही दोन्ही सूत्रे (काही वाक्य~$\varphi$ साठी) त्या
निष्पत्तीमध्ये~$\psi$ च्या आधी आलेली असली पाहिजेत.

स्वयंसिद्धकीय निष्पत्त्यांवर आधारित $\Proves$ संबंधाची व्याख्या
पुढीलप्रमाणे केली जाते: $\Gamma \Proves \varphi$ तेव्हा आणि केवळ तेव्हाच, जेव्हा
अशी निष्पत्ती असते की तिच्या शेवटच्या सूत्राच्या जागी
वाक्य~$\varphi$ असते (आणि वरच्या~(2) मुळे समर्थन मिळालेल्या त्या
निष्पत्तीमधील वाक्यांचा संच म्हणजे $\Gamma$ असे घेतले जाते).
~$\varphi$ ची अशी निष्पत्ती असेल की तिच्यासाठी~$\Gamma$ रिक्त असेल, म्हणजेच
त्या निष्पत्तीमधील प्रत्येक वाक्याला (1) किंवा~(3) मुळे समर्थन
मिळाले असेल, तर~$\varphi$ हे प्रमेय असते. उदाहरणार्थ, $\Proves
\varphi  \lif (\psi \lif (\psi \lor \varphi))$ हे दाखवणारी निष्पत्ती पुढे दिली आहे:
\begin{derivation}
  1. & $\psi \lif (\psi \lor \varphi)$ \\
  2. & $(\psi \lif (\psi \lor \varphi)) \lif (\varphi  \lif (\psi \lif (\psi \lor \varphi)))$\\
  3. & $\varphi  \lif (\psi \lif (\psi \lor \varphi))$
\end{derivation}
ओळ~1 वरील वाक्य हे $\varphi \lif (\varphi \lor \psi)$ या स्वयंसिद्धकाच्या
रूपाचे आहे ($\varphi$ आणि $\psi$ यांच्या भूमिका उलट आहेत). ओळ~2 वरील वाक्य
$\varphi \lif (\psi \lif \varphi)$ या स्वयंसिद्धकाच्या रूपाचे आहे. त्यामुळे दोन्ही
ओळींना समर्थन मिळते. ओळ~3 ला विध्यात्मक अनुमानामुळे समर्थन मिळते: त्याचा
संक्षेप $\theta$ असा केल्यास, ओळ~2 चे रूप $\chi \lif \theta$ असे आहे, आणि त्यातील
$\chi$ म्हणजे $\psi \lif (\psi \lor \varphi)$, म्हणजेच ओळ~1 होय.

$\Gamma \Proves \lfalse$ असेल, तर संच $\Gamma$ विसंगत असतो. संपूर्ण
स्वयंसिद्धकीय पद्धत कोणत्याही~$\varphi$ साठी $\lfalse \lif \varphi$ हेही सिद्ध करेल;
त्यामुळे $\Gamma$ विसंगत असेल, तर कोणत्याही~$\varphi$ साठी $\Gamma \Proves \varphi$.

तर्कशास्त्रासाठी स्वयंसिद्धकीय निष्पत्त्यांच्या पद्धती सर्वप्रथम गोटलोप
फ्रेग यांनी त्यांच्या १८७९ मधील \emph{Begriffsschrift} या ग्रंथात दिल्या;
म्हणून हा ग्रंथ आधुनिक तर्कशास्त्रातील पहिला ग्रंथ मानला जातो. ॲल्फ्रेड
नॉर्थ व्हाइटहेड आणि बर्ट्रंड रसेल यांच्या \emph{Principia Mathematica} मध्ये,
तसेच १९२० च्या दशकात डाव्हीट हिल्बर्ट आणि त्यांच्या विद्यार्थ्यांकडून, या
पद्धतींना परिपूर्ण रूप मिळाले. म्हणून त्यांना अनेकदा ``फ्रेग पद्धती'' किंवा
``हिल्बर्ट पद्धती'' म्हणतात. एखाद्या तर्कशास्त्रासाठी स्वयंसिद्धकीय पद्धत
शोधणे बहुधा सोपे असल्यामुळे या पद्धती अतिशय बहुपयोगी आहेत. निष्पत्त्यांची रचना अत्यंत सोपी असून त्यांत फक्त एक किंवा दोन निगमन नियम असतात;
त्यामुळे \emph{त्यांच्याविषयीच्या} गोष्टी सिद्ध करणेही तुलनेने सोपे असते. तथापि,
प्रत्यक्ष व्यवहारात त्यांचा वापर फार कठीण असतो; म्हणजेच सिद्धता शोधणे आणि
लिहिणे कठीण असते.



